\documentclass[a4paper]{article}

%% Language and font encodings
\usepackage[english]{babel}
\usepackage[T1]{fontenc}

%% Sets page size and margins
\usepackage[top=3cm,bottom=2cm,left=3cm,right=3cm,marginparwidth=1.75cm]{geometry}

%% Useful packages
\usepackage{amsmath}
\usepackage{amssymb}
\usepackage{graphicx}
\usepackage[colorinlistoftodos]{todonotes}
\usepackage[colorlinks=true, allcolors=blue]{hyperref}


\title{Effects of Underflow on Solving Linear Systems (Excerpt)}
\author{James Demmel\\ Computer Science Division \\ Electronic Research Laboratory \\
  U.C. Berkeley}
\date{May 1981}


\begin{document}

\maketitle

\begin{abstract}
  \url{https://www2.eecs.berkeley.edu/Pubs/TechRpts/1983/CSD-83-128.pdf}\newline
  \emph{This is an excerpt, prepared by Andrew W. Appel in 2025, of some parts of the
  tech. report relevant to accuracy bounds for forward substitution and back substitution.}
  \newline
  In this paper we examine the effects of underflow on solving systems of linear equations using Gaussian Elimination.  Our goal is to decide if reliable software for solving linear
  systems in the presence of underflow can be written at reasonable cost.
  We constrast the utilities of gradual underflow and ``store zero'',
  and show that only by using gradual underflow can we achieve reliability easily.
  \end{abstract}

\subsection*{1. Introduction}

~.~.~.

\subsection*{2. Examples}

~.~.~.

Let $\epsilon$ be the rounding error of the arithmetic, and $\lambda$ be the
underflow threshold.  Suppose, for example, a binary floating point number is
represented as $f\cdot 2^e$ where $f$ lies between 1 and 2 with an $n$ bit fraction,
and $e$ is an integer exponent in the range $e_\mathrm{min} \le e \le e_\mathrm{max}$.
Then $\epsilon=2^{-n}$ (the difference between 1 and the next largest number).
The underflow threshold is $\lambda=2^{e_\mathrm{min}}$; this is the smallest \ldots
normalized number for G.U. \ldots Denormalized numbers lie between 0 and $2^{e_\mathrm{min}}$
in G.U. arithmetic for which the smallest nonzero number is thus $\mu = \epsilon \lambda$.
We will give a more precise model of rounding and underflow errors in the next section.

\subsection*{3. Error Analysis}

~.~.~.

\newcommand\fl[1]{\ensuremath{\mathit{fl}(#1)}}

Let $\blacksquare$ be one of the operations $(+,-,\ast,/)$ and let
$\fl{a \blacksquare b}$ denote the floating point result of the indicated
computation. \ldots To take underflow into account, we write ([Kahan and Palmer, 1979])
\begin{equation}\tag{11}
\fl{a \blacksquare b} = (a \blacksquare b)(1 + e)+\eta\qquad
\mathrm{unless}~a\blacksquare b ~\mathrm{overflows.}
\end{equation}
In the case of G.U. we have the following constraints on $e$ and $\eta$:
\begin{enumerate}
\item[(1)] $|e|\le \epsilon ~\mathrm{and}~ |\eta|\le \lambda \epsilon$,
\item[(2)] $\eta \times e = 0$
\item[(3)] $\eta=0$ if $\blacksquare$ is either addition or subtraction.
\end{enumerate}

~.~.~.

We define the function $\mathit{UN}(x)$ to be
\[
UN(x)=\left\lbrace\begin{array}{ll}1 & \mathrm{if}~x~\mathrm{underflows}~(x<\lambda)\\
    0 & \mathrm{if~it~does~not}\end{array}\right.
  \]

We assume no overflow occurs.

We ignore terms which are $O(\epsilon^2)$ or $O(\epsilon\mu)$.  This allows us
to make the following substitutions when convenient:
\begin{itemize}
\item[] Replace $1/(1+e_1)$ by $1+e_2$,
\item[] Replace $(1+e_1)^j$ by $1 + je_2$,
\item[] Replace $\prod_{j=1}^m(1+e_j)/\prod_{i=1}^n (1+e'_i)$ by $1+(n+m)e$
\item[] Replace $\eta_1 \ast (1+e_1)^j$ by $\eta_2$.
\end{itemize}

By replacing every appearance of an original datum $a$ by $1. \cdot a$,
and using the above formula for error in multiplication, we can take into
account effects of rounding and underflow errors in the original data.\footnote{Editors note: Ugh.  It's provable that $1. \cdot a=a$ exactly, so this notational trick does not
appeal to me.}

\newcommand{\norm}[1]{\ensuremath \lVert #1 \rVert}
$\norm{A}$ and $\norm{b}$ denote the norms of matrix $A$ and vector $b$.
$\norm{A}_\infty$ denotes the infinity norm of $A$, and similarly for
$\norm{b}_\infty$.  $k(A)=\norm{A}_\infty\norm{A^{-1}}_\infty$ denotes the condition number
of the matrix $A$.

$|A|~(|b|)$ denotes the matrix (vector) whose entries are the absolute values
  of the corresponding entries in $A$ ($b$).  Inequalities like $|A|>|B|$
  $(|a|>|b|)$ are meant componentwise.

\subsection*{3.2 Approach}

As stated in the introduction, we use backward error analysis.  Thus, when trying
to solve
\begin{equation}\tag{12}
  Ax = b
\end{equation}
for $x$ we obtain an approximation $\hat{x}=x~\delta x$ which satisfies the
perturbed problem
\begin{equation}\tag{13}
  (A+\delta A)\hat{x}=b+\delta b~.
\end{equation}
The task of backwards error analysis is to obtain bounds on $\delta A$ and $\delta b$.
These bounds can be used in turn to bound the residual $r$:
\begin{equation}\tag{14}
  r = A \hat{x}-b = -\delta A \hat{x} + \delta b
\end{equation}

~.~.~.


\subsection*{3.3 Error Analysis Details}

~.~.~.

\paragraph*{Proposition 2.}~

\newcommand\float[1]{\ensuremath \mathit{float}(#1)}
To analyze solving a triangular system, we need another expression for the error in
computing an inner product:
\begin{equation}\tag{31}
  \float{\sum_{i=1}^n a_i b_i} = \sum_{i=1}^n a_i b_i (1 + E_i) + \eta ~.
\end{equation}
In the absence of underflow we have
\begin{equation}\tag{32}
\begin{array}{l}
  |E_1| \le n \epsilon \\
  |E_j| \le (n+2-j)\epsilon \quad \mathrm{if}\quad j>1~.
\end{array}
\end{equation}

In the case of G.U.\footnote{Editor's note: G.U. = Gradual Underflow, that is,
a floating-point format that supports subnormal numbers.  The technical report also covers the case of ``S.Z.'' or ``store zero'', which we omit here in all lemmas and theorems.}
we have the same bounds on the $|E_i|$, and
\begin{equation}\tag{34}
|\eta| \le n \lambda \epsilon
\end{equation}

\paragraph*{Proof.} Letting
\begin{eqnarray*}
  S_m & = & \fl{\sum_{i=1}^m a_i b_i} \\
  & = & \fl{S_{m-1} + a_m b_m} \\
  & = & (S_{m-1} + a_m b_m(1+e_m)+\eta_m)(1+w_{m-1})+z_{m-1}
\end{eqnarray*}
It is easy to prove by induction on $m$ that
\begin{eqnarray*}
  S_m & = & a_1b_1(1+e_1)\prod_{i=1}^{m-1}(1+w_i)+\eta_1 \prod_{i=1}^{m-1}(1+w_i)+\\
  & & \sum_{k=2}^m a_k b_k (1+e_k)\prod_{i=k-1}^{m-1}(1+w_i)+
  \sum_{k=2}^m \eta_k \prod_{i=k-1}^{m-1}(1+w_i) +\\
  & & \sum_{k=1}^{m-1}z_k \prod_{i=k+1}^{m-1}(1+w_i).
\end{eqnarray*}
$z_{m-1}=0$ always for G.U. \ldots Thus \ldots $z_{m-1}\eta_m = 0$, so the
contribution from underflow is at most $n$ times the bound for $\eta_m$
and $z_{m-1}$.  Approximating $(1+e_k)\prod_{i=r}^{n-1}(1+w_i)$ by
$1+(n-r+1)\epsilon$ we get the result.  Note that the $\sum_{i=1}^n a_i b_i(1+E_i)$
term is the same term as would have been given by round-off alone.  Q.E.D.

~.~.~.

\subsection*{Lemma 4: Error Analysis of Solving a Triangular System}

Consider the lower triangular system $Ly=b$ where $L$ does not necessarily
have ones on the diagonal.  (The analyses when only ones are on the
diagonal or L is upper triangular are similar.)  Then in solving
$Ly=b$ using the program in Appendix 2 (where inner products are
accumulated with roundoff error $\epsilon_\mathrm{s}$ and
smallest number $\mu_\mathrm{s}$ we obtain $\hat{y}=y+\delta y$ which
actually satisfies $(L+\delta L)\hat{y}= b + \delta b$, where
\begin{equation}\tag{80}
  |\delta L_{ij}| \le |L_{ij}| \cdot \left\lbrace
  \begin{array}{lll}
    \epsilon & \mathrm{if} & i=j=1 \\
    (i-1)\epsilon_\mathrm{s} & \mathrm{if} & i>j=1 \\
    (i+1-j)\epsilon_\mathrm{s} & \mathrm{if} & i>j>1 \\
    \epsilon+\epsilon_\mathrm{s} & \mathrm{if} & i=j>1 
  \end{array}\right.
\end{equation}
independent of underflow.

In the absence of underflow
\begin{equation}\tag{81}
  \delta b = 0~.
\end{equation}
Using G.U.
\begin{equation}\tag{82}
  \begin{array}{lcl@{~}ll}
    \mathrm{if} & i=1 & |\delta b_i |~ \le &\lambda\epsilon|L_{11}|\mathit{UN}(y_1)
    & \mathrm{if}~y_1~\mathrm{underflows} \\
    \mathrm{if} & i>1 & |\delta b_i |~ \le &i \cdot \lambda_\mathrm{s}\epsilon_\mathrm{s}~ + & \mathrm{intermediate~underflows} \\
    & & & \lambda \epsilon |L_{ii}| \mathit{UN}(y_i) 
    & \mathrm{if}~y_i~\mathrm{underflows}
  \end{array} ~.
\end{equation}
If $L_{ii}$ is known to be 1, the only changes in the above bounds
(besides substituting 1 for $L_{ii}$) are that $E_{11} = \delta b_1 = 0$
independent of underflow.

\paragraph*{Proof:}  We apply Proposition 2 to the program in Appendix 2.
We have two cases.

\noindent Case $i=1$: In this case,
\begin{equation}\tag{84}
  y_1 = \float{b_1/L_{11}} = \frac{b_1}{L_{11}(1+E_{11})}+\nu_1
\end{equation}
so
\begin{equation}\tag{85}
  L_{11}(1+E_{11})y_1 = b_1 + \nu_1 L_{11}(1+E_{11})\mathit{UN}(y_1)
\end{equation}
where
\begin{equation}\tag{86}
  |E_{11}|\le \epsilon \quad \mathrm{and} \quad |\nu_1|\le \lambda ~\mathrm{or}~\lambda/\epsilon
\end{equation}

\noindent Case $i>1$: Now
\begin{equation}\tag{87}
\begin{array}{lcl}
  y_i & = & \float{(b_i - \sum_{j=1}^{i-1}L_{ij}y_j)/L_{ii}}\\
  & = & [(1+e_i)(b_i - \sum_{j=1}^{i-1}L_{ij}(1+E_{ij})y_j-\rho_i)+\sigma_i]/
  (L_{ii}(1+\zeta_i))+\nu_i
  \end{array}
\end{equation}
where
\begin{equation}\tag{88a}
  |E_{ij}|\le\left\lbrace \begin{array}{ll}(i-1)\epsilon_\mathrm{s} & \mathrm{if}~~j=1\\
    (i+1-j)\epsilon_\mathrm{s} & \mathrm{if}~~j>1 \end{array}\right.
\end{equation}
\begin{equation}\tag{88b}
|e_i|\le \epsilon_\mathrm{s}
\end{equation}
\begin{equation}\tag{88c}
|\zeta_i|\le \epsilon
\end{equation}
and
\begin{equation}\tag{89a}
|\rho_i|\le (i-1)\lambda_\mathrm{s}\epsilon_\mathrm{s}\qquad\mathrm{using~G.U.}
\end{equation}
\begin{equation}\tag{89b}
|\sigma_i|\le \lambda_\mathrm{s}\epsilon_\mathrm{s}\qquad\mathrm{using~G.U.}
\end{equation}
\begin{equation}\tag{89c}
|\nu_i|\le \lambda\epsilon\qquad\mathrm{using~G.U.}
\end{equation}

In the absence of underflow, $\rho_i=\sigma_i=\nu_i=0$.
Rearranging equation (87), we obtain
\begin{equation}\tag{90}
  L_{ii}(1+\zeta_i)y_i = (1+e_i)(b_i-\sum_{j=1}^{i-1}L_{ij}(1+E_{ij})y_j - \rho_i)
  + \sigma_i + \nu_iL_{ii}(1+\zeta_i)
\end{equation}
or
\begin{equation}\tag{91}
  b_i = \sum_{j=1}^{i-1}L_{ij}(1+E_{ij})y_j+L_{ii}\left(\frac{1+\zeta_i}{1+e_i}\right)y_i
  + \rho_i - \frac{\sigma_i}{1+e_i} - \nu_i L_{ii} \left(\frac{1+\zeta_i}{1+e_i}\right)~~.
\end{equation}
Using the bounds in (88) and (89) in (91), we obtain the bounds
in the statement of the lemma.  Q.E.D.

~.~.~.

\subsection*{Appendix 2: Program Listing of Solution of Both Triangular Systems Following
  Triangular Decomposition}

\begin{tabbing}
array B[1..N] of real;\\
array X[1..N] of real;\\
array Y[1..N] of real;\\
/* Again, B,X, and Y could occupy the same locations in memory. Arrays\\
A,L and U are defined above.  As in Appendix 1, precisions of all computations\\
except entries of B, X, and Y are shown. */\\
\\
/* Forward Substitution: $LY=B$ */\\
\\
for i=1 to N do\\
~~~~~\=begin\\
\>SUM=0;\\
\>for k=1 to i-1 do SUM=SUM+L[i,k]*Y[k]; /* precision $(\epsilon_\mathrm{s},\mu_\mathrm{s})$ */\\
\>Y[i]=B[i]-SUM; /* precision $(\epsilon_\mathrm{s},\mu_\mathrm{s})$\\
\>/* \=Editor's note: After Cholesky factorization, where L[i,i] is not 1, must\\
\>\>also divide by L[i,i] with precision $(\epsilon,\mu)$ */\\
\\
/* Back Substitution: UX=Y */\\
for i=N downto 1 do\\
\>begin\\
\>SUM=0;\\
\>for k=i+1 to N do SUM = SUM + U[i,k]*X[k]; /* precision $(\epsilon_\mathrm{s},\mu_\mathrm{s})$ */\\
\>X[i]=(Y[i]-SUM)/U[i,i]; /* precision  $(\epsilon_\mathrm{s},\mu_\mathrm{s})$ */\\
/* Editor's note: in the line above, the division is done at precision $(\epsilon,\mu)$ */\\
\end{tabbing}

\newpage

\section*{Application to libValidSDP}

That concludes the excerpt of Demmel's 1981 technical report.  Now we
explain what's proved in libValidSDP.

Instead of Demmel's Proposition 2, we will be using
the following lemma:

\paragraph*{Rump/Jeannerod Lemma 2.1:} Error analysis of $(c-\sum a_i b_i)/b_k$.
  (libvalidsdp.cholesky.lemma\_2\_1)
  \newline
  Let $\hat{y} = \fl{(c-\sum_{0\le i<k} a_i b_i)/b_k}$ computed by repeated subtraction left to right.  Assume $b_k\not= 0$. Then
\begin{equation}\tag{lemma2.1}
  |b_k \hat{y} - (c-\sum_{0 \le i<k}a_i b_i)| < (k+1)\epsilon |b_k\hat{y}|\sum_{0\le i < k}|a_ib_i|
  + (1+(k+1)\epsilon)(k+1+|b_k|)\lambda\epsilon.
  \end{equation}
  
\subsection*{Error Analysis of Solving a Triangular System}

\paragraph*{Demmel Lemma 4, zero-based indexing.} Consider the lower triangular system $Ly=b$ where $L$ does not necessarily
have ones on the diagonal.  (The analyses when only ones are on the
diagonal or L is upper triangular are similar.)  Then in solving
$Ly=b$ using the program in Appendix 0.2 (where inner products are
accumulated with roundoff error $\epsilon_\mathrm{s}$ and
smallest number $\mu_\mathrm{s}$ we obtain $\hat{y}=y+\delta y$ which
actually satisfies $(L+\delta L)\hat{y}= b + \delta b$, where
\begin{equation}\tag{0.80}
  |\delta L_{ij}| \le |L_{ij}| \cdot \left\lbrace
  \begin{array}{lll}
    \epsilon & \mathrm{if} & i=j=0 \\
    i\epsilon_\mathrm{s} & \mathrm{if} & i>j=0 \\
    (i+1-j)\epsilon_\mathrm{s} & \mathrm{if} & i>j>0 \\
    \epsilon+\epsilon_\mathrm{s} & \mathrm{if} & i=j>0 
  \end{array}\right.
\end{equation}
independent of underflow.

However, in all these cases we will prove a
weaker bound because that allows us to
re-use the Rump/Jeannerod Lemma 2.1 for $c - \sum x_i y_i$.

\paragraph*{Theorem Demmel\_L4\_adapted:} 
\begin{equation}\tag{demmel\_lemma\_4}
  |\delta L_{ij}| \le |L_{ij}|n\epsilon\qquad\mathrm{and}
\end{equation}
\begin{equation}
  \begin{array}{lcl@{~}ll}
    \mathrm{if} & i=0 & |\delta b_i |~ \le &\lambda\epsilon|L_{00}|\mathit{UN}(y_0)
    & \mathrm{if}~y_0~\mathrm{underflows} \\
    \mathrm{if} & i>0 & |\delta b_i |~ \le &(1+(i+1)\epsilon)(i+1) \cdot \lambda_\mathrm{s}\epsilon_\mathrm{s}~ + & \mathrm{intermediate~underflows} \\
    & & & (1+(i+1)\epsilon)\lambda \epsilon |L_{ii}| \mathit{UN}(y_0) 
    & \mathrm{if}~y_i~\mathrm{underflows}
  \end{array} ~.
\end{equation}
\noindent The factor $(1+(i+1)\epsilon)$ was omitted from Demmel's original but is required for correctness. 

\paragraph*{Proof.}  From libValidSDP lemma\_2\_1\_aux, by a derivation of several lines.  Q.E.D.

\subsection*{Solving a linear system by Cholesky+Substitution}

Let $\hat{x}$ be the result of solving $Ax=b$ by Cholesky decomposition $A=R^T R$ followed by forward substitution $R^Ty=b$ then back substitution $Rx=y$.

\paragraph*{Lemma th\_2\_3\_aux3: Accuracy of Cholesky decomposition.}
Suppose $A$ is an $n\times n$ symmetric positive definite floating-point matrix,
with $(n+1)\epsilon < 1$.
Let $\hat{R}$ be the result of floating-point Cholesky decomposition (inner-product method)
of  $A$ into $\hat{R}^T \hat{R}$.
Let $\delta A = A - \hat{R}^T \hat{R}$.
Let the vector $R_j$ be the $j$th column of $R$.
Let $a_\mathrm{max}$ be the maximum diagonal element of $|A|$.
Then
\begin{equation}\tag{2.3aux3}
  |\delta A _{ij}| < (\mathrm{min}(i,j)+2)\epsilon \norm{R_i}_2 \norm{R_j}_2
  + 4 \lambda\epsilon (n+1+a_\mathrm{max}).
\end{equation}

\paragraph*{Proof.} See Rump+Jeannerod, or see libValidSDP.  Q.E.D.

\paragraph*{Remark.}  Let $A$ and $\hat{R}$ be as above.
Then
\begin{equation}\tag{2.3aux1}
  \norm{\hat{R}_j}_2 \le \sqrt{(A_{jj}+2j\lambda\epsilon)/(1-(j+2)\epsilon)}.
  \end{equation}

\paragraph{Proof.} LibValidSDP   th\_2\_3\_aux1.


\paragraph{Corollary.} Combining (2.3aux1) and (2.3aux2), we have,
\begin{equation}\tag{ch\_back\_err}
\begin{array}{rl}
|\delta A_{ij}| & <
(\mathrm{min}(i,j)+2)\epsilon
\sqrt{(A_{ii}+2i\lambda\epsilon)/(1-(i+2)\epsilon)}
\sqrt{(A_{jj}+2j\lambda\epsilon)/(1-(j+2)\epsilon)}
+ 4 \lambda\epsilon (n+1+a_\mathrm{max}) \\
& \le
(n+1)\epsilon
\sqrt{(a_\mathrm{max}+2(n-1)\lambda\epsilon)/(1-(n+1)\epsilon)} \\
& \qquad \cdot
\sqrt{(a_\mathrm{max}+2(n-1)\lambda\epsilon)/(1-(n+1)\epsilon)}
+ 4 \lambda\epsilon (n+1+a_\mathrm{max}) \\
& =
(n+1)\epsilon
(a_\mathrm{max}+2(n-1)\lambda\epsilon)/(1-(n+1)\epsilon)
+ 4 \lambda\epsilon (n+1+a_\mathrm{max}) \\
\end{array}
\end{equation}

\paragraph*{Lemma.} We obtain bounds on $\delta R^T \hat{R}$ (et cetera) by
combining (2.3aux1) and (demmel\_lemma\_4):
\begin{equation}\tag{csolve\_back\_err\_aux1}
\begin{array}{rl}
|(\delta R^T\hat{R})_{ij}| & \le (|\delta R|^T|\hat{R}|)_{ij}| \\
& \le  \sum_k (k+1)\epsilon |\hat{R}_{ki}|\cdot|\hat{R}_{kj}| \qquad \mathrm{(by~Lemma~0.4)}\\
& \le n\epsilon \norm{\hat{R}_i}_2 \norm{\hat{R}_j}_2 \qquad \mathrm{(by~Cauchy-Schwartz)}\\
& \le n\epsilon \sqrt{(A_{ii}+2i\lambda\epsilon)/(1-(i+2)\epsilon)} \sqrt{(A_{jj}+2j\lambda\epsilon)/(1-(j+2)\epsilon)} \qquad\mathrm{(2.3aux1)}\\
& \le n\epsilon \sqrt{(a_\mathrm{max}+2(n-1)\lambda\epsilon)/(1-(n+1)\epsilon)} \sqrt{(a_\mathrm{max}+2(n-1)\lambda\epsilon)/(1-(n+1)\epsilon)} \\
& = n\epsilon (a_\mathrm{max}+2(n-1)\lambda\epsilon)/(1-(n+1)\epsilon) \\
\end{array}
\end{equation}

\paragraph{Lemma.} Bounds on $\delta R^T \delta R$ as follows:
\begin{equation}\tag{csolve\_back\_err\_aux2}
  \begin{array}{rl}
    |(\delta R ^T \delta R)_{ij}| & \le (|\delta R^T| \cdot |\delta R|)_{ij} \\
    & \le \sum_k |\delta R^T|_{ik} |\delta R|_{kj} \\
    & = \sum_k |delta R|_{ki} |\delta R|_{kj} \\
    & \le \sum_k (k+1)^2\epsilon^2 |\hat{R}_{ki}|\cdot|\hat{R}_{kj}| \\
    & \le n^2\epsilon^2 \norm{\hat{R}_i}_2 \norm{\hat{R}_j}_2 \\
    & \le n^2\epsilon^2 (a_\mathrm{max}+2(n-1)\lambda\epsilon)/(1-(n+1)\epsilon).
  \end{array}
\end{equation}


\paragraph*{Theorem: Cholesky factor + substitution.}
Suppose $A$ is an $n\times n$ symmetric positive definite floating-point matrix
with minimum eigenvalue $>\sqrt{\lambda}$, with $(n+1)\epsilon < 1$.
Suppose $b$ is a floating-point vector of length $n$.
Suppose $\hat{x}$ is the result of floating-point Cholesky factorization of
$A$ into $\hat{R}^T\hat{R}$ followed by forward substitution finding $\hat{y}$ such
that $\hat{R}^T \hat{y}=b$ and then $\hat{R}\hat{x}=\hat{y}$.
Then
\begin{equation}\tag{xxx}
  \mathrm{fill~this~in~\ldots}
\end{equation}

\paragraph*{Proof.}
If it were not for the underflow terms $\delta b$ and $\delta y$, we would have
\begin{equation}\tag{csolve\_back\_err}
  \begin{array}{rl}
    b &= (\hat{R}^T+\delta R^T)\hat{y}\\
    & = (\hat{R}^T+\delta R^T)(\hat{R}+\delta R)\hat{x} \\
    & = (\hat{R}^T\hat{R}+\delta R^T\hat{R} +\hat{R}^T\delta R  +\delta R^T \delta R)\hat{x} \\
    & = (A + \delta A +\delta R^T\hat{R} +\hat{R}^T\delta R  +\delta R^T \delta R)\hat{x} \\
    & = A\hat{x} + (\delta A +\delta R^T\hat{R} +\hat{R}^T\delta R  +\delta R^T \delta R)\hat{x}. \\
  \end{array}
\end{equation}

From (2.3aux3) we have
$  |\delta A _{ij}| < (\mathrm{min}(i,j)+2)\epsilon \norm{R_i}_2 \norm{R_j}_2
+ 4 \lambda\epsilon (n+1+a_\mathrm{max})$.

From 


\subsection*{Appendix 0.2: Program Listing of Solution of Both Triangular Systems Following
  Triangular Decomposition}

\begin{tabbing}
array B[N] of real;\\
array X[N] of real;\\
array Y[N] of real;\\
/* Again, B,X, and Y could occupy the same locations in memory. Arrays\\
A,L and U are defined above.  As in Appendix 1, precisions of all computations\\
except entries of B, X, and Y are shown. */\\
\\
/* Forward Substitution: $LY=B$ */\\
\\
for (i=0; i<N; i++)\\
~~~~~\=begin\\
\>SUM=0;\\
\>for (k=i; k<i-1; k++) do SUM=SUM+L[i,k]*Y[k]; /* precision $(\epsilon_\mathrm{s},\mu_\mathrm{s})$ */\\
\>Y[i]=B[i]-SUM; /* precision $(\epsilon_\mathrm{s},\mu_\mathrm{s})$\\
\>/* \=Editor's note: After Cholesky factorization, where L[i,i] is not 1, must\\
\>\>also divide by L[i,i] with precision $(\epsilon,\mu)$ */\\
\\
/* Back Substitution: UX=Y */\\
for (i=N-1; i>=0; i--) do\\
\>begin\\
\>SUM=0;\\
\>for (k=i+1; i<N; k++) do SUM = SUM + U[i,k]*X[k]; /* precision $(\epsilon_\mathrm{s},\mu_\mathrm{s})$ */\\
\>X[i]=(Y[i]-SUM)/U[i,i]; /* precision  $(\epsilon_\mathrm{s},\mu_\mathrm{s})$ */\\
/* Editor's note: in the line above, the division is done at precision $(\epsilon,\mu)$ */\\
\end{tabbing}
\end{document}
